% =============================================================================
% Section: System architecture  (plan item A1)
% Intended position: new Section III, before "Mathematical model".
% Notation introduced here and used later:
%   \hat e_j              runtime estimate used for job j  (e_j^raw := \hat e_j)
%   b_m(t)                instances of pool m held by running jobs at time t
%   (\hat m_j, \hat s_j)  planned pool and planned start of job j
% \revnote{...} marks items the authors must verify before submission.
% =============================================================================
\section{System Architecture}
\label{sec:architecture}

MOSPACE is a \emph{central scheduler}. The federation is a set of resource
pools $M$. An on-premises pool is a group of identical nodes at one site. A
cloud pool is one VM type of one provider. MOSPACE is the only scheduling
authority over all pools. Every job enters its queue, and it alone decides in
which pool each job runs and when the job starts. The pools only execute: a
pool starts a job on whole nodes when MOSPACE tells it to and reports when
the job starts and ends.

This assignment of authority differs from existing multi-cluster
mechanisms. Slurm's multi-cluster mode and Slurm federation choose a cluster
for each job and leave the start time to that cluster's
scheduler~\cite{Yoo2003,Jette2023}. Two-level placement~\cite{Wang2025Efficient}
and cloud-bursting policies~\cite{Luo24,ambati2021good} decide one job at a
time. With paid capacity, decisions about different jobs are coupled: a free
on-premises slot is worth most to the competing job whose cloud alternative
is the most expensive. Only a scheduler that plans all pending jobs together
can make this trade-off, so MOSPACE keeps one plan for the whole queue. This
corresponds to an operator that runs several node groups or sites and holds
its own cloud accounts, which is the setting of
Section~\ref{sec:experimental_setup}. Table~\ref{tab:arch-req} lists what the
design needs from the platform and how the evaluation provides it.

\begin{figure}[t]
  \centering
  % =============================================================================
% MOSPACE architecture (TikZ). Single column (3.5in); \input inside figure.
% Colours and pool macros: figures/architecture_preamble.tex.
% =============================================================================
\begin{tikzpicture}[
  font=\scriptsize,
  >={Stealth[length=1.9mm,width=1.4mm]},
  comp/.style={draw=archaccent, fill=white, rounded corners=1.5pt, line width=0.6pt,
               minimum width=2.45cm, text width=2.35cm, minimum height=0.95cm,
               align=center, inner sep=1.5pt},
  flow/.style={->, line width=0.6pt, draw=archaccent},
  ctrl/.style={->, line width=0.7pt},
  ev/.style={->, line width=0.6pt, draw=black!65, dash pattern=on 2.2pt off 1.4pt},
  grp/.style={rounded corners=2pt, line width=0.6pt, inner sep=0pt},
]
% ---------------------------------------------------------------- scheduler
\node[comp] (queue) at (1.35,0)    {\textbf{Job queue}\\ all pending jobs};
\node[comp] (plan)  at (4.3,0)     {\textbf{Planner}\\ arc-flow MILP or GA};
\node[comp] (disp)  at (7.25,0)    {\textbf{Dispatcher}\\ starts due jobs};
\node[comp] (mon)   at (4.3,-1.55) {\textbf{State monitor}\\ running jobs, free nodes};

\draw[flow] (queue) -- (plan);
\draw[flow] (plan)  -- (disp);
\draw[flow] (mon)   -- node[right, align=left] {re-plan if the\\ plan is invalid} (plan);

\begin{scope}[on background layer]
  \node[grp, draw=archaccent, fill=archbroker, fit={(0,0.95) (8.6,-2.25)}] (sched) {};
\end{scope}
\node[anchor=north east, font=\scriptsize\bfseries, text=archaccent] at (8.52,0.92)
  {MOSPACE central scheduler};

% submission
\node[draw, rounded corners=1.5pt, line width=0.6pt, minimum width=2.45cm,
      minimum height=0.55cm, font=\scriptsize\bfseries] (users) at (1.35,1.75) {Users};
\draw[ctrl] (users) -- node[right, pos=0.22] {submit} (queue);

% ----------------------------------------------------------- execution layer
\node[grp, draw=black!65, fill=archsite, minimum width=4.3cm, minimum height=1.9cm]
  (onp) at (2.15,-4.3) {};
\node[anchor=north west, font=\scriptsize\bfseries] at ([shift={(0.06,-0.03)}]onp.north west)
  {On-premises pools};
\archpool{0.15}{-3.95}{4}{$m_1$ (site A)}
\archpool{0.15}{-4.35}{5}{$m_2$ (site A)}
\archpool{0.15}{-4.75}{2}{$m_3$ (site B)}
\node[anchor=south east, font=\scriptsize\itshape, align=right] at ([shift={(-0.06,0.05)}]onp.south east)
  {fixed capacity, no charge};

\node[grp, draw=archcloudln, fill=archcloud, minimum width=4.0cm, minimum height=1.9cm]
  (cl) at (6.6,-4.3) {};
\node[anchor=north west, font=\scriptsize\bfseries] at ([shift={(0.06,-0.03)}]cl.north west)
  {Cloud pools};
\archvm{4.75}{-3.95}{2}{$m_4$ (AWS type)}
\archvm{4.75}{-4.35}{1}{$m_5$ (Azure type)}
\node[anchor=south east, font=\scriptsize\itshape, align=right] at ([shift={(-0.06,0.05)}]cl.south east)
  {on demand, priced};

\begin{scope}[on background layer]
  \node[grp, draw=black!40, dash pattern=on 3pt off 1.5pt, fit={(-0.1,-3.15) (8.7,-5.35)}] (exec) {};
\end{scope}

% decisions down, events up
\draw[ctrl] (disp.south) -- node[left, align=right] {start $j$ on\\ pool $\hat m_j$} (7.25,-3.15);
\draw[ev] (4.45,-3.15) -- node[right] {start / end events} (4.45,-2.025);
\end{tikzpicture}

  \caption{MOSPACE architecture. The central scheduler plans every pending
  job and starts each one on its planned pool $\hat m_j$. The pools only
  execute jobs and report events. Shaded squares are busy nodes; dashed
  squares are cloud VMs that are not running.}
  \label{fig:architecture}
\end{figure}

\subsection{Components}
\label{sec:arch-components}

Fig.~\ref{fig:architecture} shows the four components of the scheduler.

\emph{Job queue.} The queue holds every pending job $j$ with its arrival
time $a_j$, demands $q_j^r$, I/O volume $\iota_j$, priority $w_j$ and runtime
estimate $\hat e_j$. The model uses $e^{\mathrm{raw}}_j=\hat e_j$.

\emph{State monitor.} The monitor tracks running jobs and free nodes from the
start and end events of the pools. A running job holds its nodes until its
estimated end, which gives the committed load $b_m(t)$. Only $N_m-b_m(t)$
instances of pool $m$ are free for planning at time $t$.

\emph{Planner.} The planner builds the model of Section~\ref{sec:model} from
the queue and the committed load and solves it with the arc-flow MILP or the
GA of Section~\ref{sec:solutions}. The result is a plan: a pool $\hat m_j$ and a
start $\hat s_j$ for every pending job. Identical nodes of an on-premises
site form one pool; a job may use the pools that its partition request
allows. Each cloud VM type is its own pool. The horizon $H$ is extended until
an all-on-premises list schedule of the queue fits, so a zero-cost schedule
always exists, as Section~\ref{sec:model} assumes.

\emph{Dispatcher.} At time $\tau$, the dispatcher starts every job with
$\hat s_j\le\tau$ on pool $\hat m_j$ and allocates whole nodes to it. A job
may also start before $\hat s_j$ on its planned pool if this delays no other
planned job. This is conservative backfilling against the plan, and it
changes neither a planned completion nor a cost.

\subsection{Decision loop}
\label{sec:arch-loop}

The scheduler runs once per event timestamp, after all submissions and job
ends at that timestamp have been applied. It re-plans the whole queue only
if the current plan is invalid, which happens when (i) a new job has
arrived, (ii) a running job has ended before or after its estimated end, or
(iii) a planned start has failed. Otherwise the dispatcher keeps following
the plan. Following the plan loses nothing when no event deviates from it,
because the part of an optimal plan that has not yet started is optimal for
the remaining jobs. Running jobs are never moved.

Jobs already pending when MOSPACE takes over a system are a backlog that the
current plan did not schedule. They are started on-premises in
first-come-first-served order before the planner places later jobs on the
remaining capacity.

\begin{table}[t]
\centering
\caption{What the central scheduler requires from the platform, and how the
evaluation provides it.}
\label{tab:arch-req}
\footnotesize
\renewcommand{\arraystretch}{1.12}
\begin{tabularx}{\columnwidth}{@{}>{\raggedright\arraybackslash}X>{\raggedright\arraybackslash}p{3.3cm}@{}}
\toprule
Requirement & In the evaluation \\
\midrule
Exclusive control: all jobs on the pools pass through the scheduler &
  All trace jobs; jobs running or pending at the start are included \\
A runtime estimate $\hat e_j$ for each job &
  Runtime recorded in the trace\revnote{or walltime / prediction, if added} \\
Start a job on whole nodes of a given pool &
  Simulator allocator, exclusive nodes \\
Job start and end events &
  Simulator events \\
Price, boot time and quota of each cloud pool &
  On-demand list prices; boot time and quota per experiment \\
\bottomrule
\end{tabularx}
\end{table}

\emph{Scope.} The evaluation runs the queue, planner and dispatcher logic
unchanged inside the simulator, which plays the role of the pools. It does
not include a production deployment. The architecture fixes the model's
assumptions: jobs are single-node parallel jobs with exclusive allocation,
runtimes are estimates, and the cost objective counts cloud expenditure
only.
\revnote{Report the decision time per call (p50/p95/max) and the fraction of
calls that needed a solve.}
